\section{Before you think and choose}

\begin{quotex}
Avant que tu penses et que tu choisisses, la société s'empare de ton entité et la façonne, selon son droit. Des que tu penses et choisis, efface les plis reçus, c'est-a-dire libère-toi des habitudes contemporaines, selon ton devoir.

Pour choisir, sache que tu as trois destins : tu peux être un animal comme ce décadent que les superficiels nomment sauvage; un animique comme tout le monde, un spirituel comme saint Thomas, ou Dante. Animal, sois beau; animique, sois bon; spirituel, cherche le Graal.

Pour embellir, animise tes instincts ; pour t'adoucir, spiritualise tes sentiments ; pour tendre à l'absolu, développe en toi l'abstraction.

\end{quotex}
\flright{\textsc{Joséphin Péladan}, \textit{Comment on devient Mage}, II, La Société)}

\textbf{Joséphin Péladan} was a key figure in the 19th century French occult revival. Along with Papus and De Guaita, he recreated the Rosicrucian order. Material that used to be secret — for reasons that will be obvious — were openly revealed. For the initiate, it doesn't much matter, because the unitiated will not understand it. However, the misunderstanding can have consequences, especially in an open society that knows nothing of self-regulation. 

For the initiate, the relationship with society can be problematic. Society operates under its own laws and facts, which supercede the individual and require him to be unconscious. The initiate, on the other hand, strives to become more conscious, transcending both the human state and society. In \textit{How to become a Mage}, Peladan writes:

\begin{quotex}
Baptism [as an initiation rite] makes us children of God, but Society dooms us to evil through is laws and education. Faith enlightens us, but it is in perpetual conflict with Society. The initiate, in order to make the grace of baptism full and effective, must renounce anew Society, its boundaries, its crimes, so that the fear of God makes him prefer the intimate greatness to the degrading favours of the country. 

\end{quotex}
The following passage is taken from the chapter on Society. Note how it coincides with the Three Stages of Man\footnote{\url{https://gornahoor.net/?p=214}}.

\begin{quotex}
Before you think and choose, society takes over your being and moulds it, as is its right. Once you think and choose, remove those received imprints, that is to say, liberate yourself from contemporary habits, as is your right. 

\end{quotex}
The uninitiated man takes as the norm the world as it appears to him; this is not the real world, but is a second reality, or a shadow world hiding its true source. Such a man lives in that second reality. The initiate begins to wake up and must see through that second reality in order to become fully conscious. Of course, in our day and age, everyone thinks he is a rebel and boasts about flaunting societal mores. Unable to conceive of any higher goal, this faux rebel can only act out against the taboos against sex, drugs, and other anti-social behaviour. This is just another trap, another way of being moulded. What better way to keep someone from liberating himself than to convince him he is already a ``free spirit"?

\begin{quotex}
In order to choose, know that you have three destinies. You can be:

\begin{enumerate}
\item an animal like the decadent that the superficial call savage; 
\item a soulish man, like everyone else; 
\item a spiritual person like St Thomas or Dante. 
\end{enumerate}

Your tasks:

\begin{itemize}
\item \textbf{Animal}: be beautiful; 
\item \textbf{Soulish man}: be good; 
\item \textbf{Spiritual person}: seek the Grail. 
\end{itemize}
\end{quotex}
The three choices:

\begin{enumerate}
\item One can act on instinct alone, like an animal. This is how many people would interpret ``liberation". Such a man should at least be a beautiful animal. 
\item The soulish man experiences no distance between his thoughts and feelings and his self. Whatever he thinks and feels \emph{is} his ``self" at that moment. His option is to be ``good", as his perspective is moralism. 
\item The destiny of the spiritual man, who transcends his thoughts and feelings, is to seek the Holy Grail. 
\end{enumerate}
\begin{quotex}
In order to improve, soulify your instincts; in order to make yourself meek, spiritualise your feelings; in order to reach the absolute, develop within yourself a separation. 

\end{quotex}
This summarizes his advice:

\begin{itemize}
\item Instincts need to be made self-conscious. 
\item Meek must be understood in the sense of the beatitudes: Blessed are the Meek\footnote{\url{https://gornahoor.net/?p=359}}. 
\item Learn to transcend the Mind. Develop a sense of detachment — the Observer, the Unmoved Mover, the Self. 
\end{itemize}


\flrightit{Posted on 2010-09-03 by Cologero }
